\documentclass[11pt,a4paper]{article}
\usepackage[utf8]{inputenc}
\usepackage[T1]{fontenc}
\usepackage[french]{babel}
\usepackage{amsmath,amssymb}
\usepackage[margin=2cm]{geometry}

\title{Corrections : Exercice 1, Exercice 2 (fin) et Exercice 3}
\date{}

\begin{document}
\maketitle

%% =====================================================================
\section*{Exercice 1}

\subsection*{Consigne 1 : co\^ut pour les finalistes}

\paragraph{Fraction des \'el\`eves qui redoublent.}
La moiti\'e entre en seconde A, les $\frac25$ entrent en seconde CD, donc
\[
1-\frac12-\frac25=\frac{1}{10}
\]
des \'el\`eves redoublent.

\paragraph{Nombre de finalistes.} On part de $80$ candidats.
\begin{itemize}
  \item Apr\`es la 1\textsuperscript{re} journ\'ee, il reste $\frac34$ : $80\times\frac34=60$.
  \item Apr\`es la 2\textsuperscript{e} journ\'ee, il reste $\frac13$ : $60\times\frac13=20$.
  \item Apr\`es la 3\textsuperscript{e} journ\'ee, les $\frac25$ restants se qualifient : $20\times\frac25=8$.
\end{itemize}
Il y a donc \textbf{8 finalistes}, soit $\frac25\times\frac13\times\frac34=\frac1{10}$ des candidats.

\paragraph{Gains.} Le gain est de $500$ F au premier tour, puis triple \`a chaque tour suivant :
$500$ F, $1\,500$ F, $4\,500$ F. Un finaliste qui reste jusqu'au bout gagne donc
\[
500+1\,500+4\,500=6\,500\ \text{F}.
\]
\textbf{Co\^ut pour la phase finale :} $8\times6\,500=\mathbf{52\,000\ F}$.

\emph{Remarque : si l'on compte aussi les gains des \'elimin\'es des deux premi\`eres journ\'ees
($60\times500$ F et $20\times1\,500$ F), le total vers\'e \`a tous les candidats est de
$30\,000+30\,000+36\,000=96\,000$ F. La question porte sur les finalistes.}

\subsection*{Consigne 2 : production \`a r\'ealiser}

Il faut construire $\varphi=\dfrac{1+\sqrt5}{2}$ \`a la r\`egle et au compas, v\'erifier les
propri\'et\'es donn\'ees, puis calculer $A$ et $B$.

\paragraph{Construction.}
\begin{enumerate}
  \item Tracer une droite gradu\'ee $(x'Ox)$ avec $O(0;0)$ et $I(1;0)$.
  \item \'Elever la perpendiculaire en $I$ et placer $K$ tel que $IK=2$.
  \item Tracer $[OK]$ : par Pythagore, $OK=\sqrt{1^2+2^2}=\sqrt5$.
  \item Tracer le cercle de centre $I$ et de rayon $\sqrt5$ : il coupe $[Ox)$ en $T(1+\sqrt5\,;\,0)$.
  \item Construire le milieu de $[OT]$ (m\'ediatrice) : son abscisse est $\dfrac{1+\sqrt5}{2}=\varphi$.
\end{enumerate}

\paragraph{V\'erification des propri\'et\'es.}
\begin{itemize}
  \item $\varphi^2=\dfrac{6+2\sqrt5}{4}=\dfrac{3+\sqrt5}{2}=\varphi+1$.
  \item $\dfrac1\varphi=\dfrac{2}{1+\sqrt5}=\dfrac{\sqrt5-1}{2}=\varphi-1$.
  \item $\varphi^3=\varphi(\varphi+1)=\varphi^2+\varphi=2\varphi+1$.
\end{itemize}

\paragraph{Calcul de $A$.} On utilise $\dfrac1{\varphi^2}=(\varphi-1)^2=2-\varphi$. Alors
\[
1+\frac{1}{1+\varphi}=1+\frac1{\varphi^2}=3-\varphi,
\qquad
A=\frac{1}{3-\varphi}=\frac{2}{5-\sqrt5}=\frac{2(5+\sqrt5)}{20}=\boxed{\frac{5+\sqrt5}{10}}\approx0{,}7236.
\]

\paragraph{Calcul de $B$.} Comme $1+\varphi=\varphi^2$, on a $\sqrt{1+\varphi}=\varphi$. De l'int\'erieur vers l'ext\'erieur :
\[
\sqrt{1+\varphi}=\varphi,\qquad \sqrt{1+\sqrt{1+\varphi}}=\varphi,\qquad B=\sqrt{1+\sqrt{1+\sqrt{1+\varphi}}}=\varphi.
\]
\[
\boxed{B=\varphi=\frac{1+\sqrt5}{2}}
\]

%% =====================================================================
\section*{Exercice 2 (questions c. à f.)}

On d\'efinit
\[
A=\sqrt{x+1+2\sqrt x}-\sqrt{x+1-2\sqrt x}
\]
pour $x\ge0$. Comme $x+1+2\sqrt x=(\sqrt x+1)^2$ et $x+1-2\sqrt x=(\sqrt x-1)^2$, on a
\[
A=(\sqrt x+1)-\lvert\sqrt x-1\rvert.
\]
\emph{Remarque : la d\'efinition de $A$ n'est pas visible sur la photo ; elle est reconstitu\'ee \`a partir des questions.}

\subsection*{c. Justifier que $A\ge0$}
Comme $\sqrt x\ge0$, on a $\lvert\sqrt x-1\rvert\le\sqrt x+1$. Donc
\[
A=(\sqrt x+1)-\lvert\sqrt x-1\rvert\ge0.
\]

\subsection*{d. Si $x\in[0;1]$, alors $A=2\sqrt x$}
Pour $x\in[0;1]$, on a $\sqrt x\le1$, donc $\lvert\sqrt x-1\rvert=1-\sqrt x$. Ainsi
\[
A=(\sqrt x+1)-(1-\sqrt x)=2\sqrt x.
\]

\subsection*{e. Valeur de $\sqrt{\sqrt2+2\sqrt{\sqrt2-1}}-\sqrt{\sqrt2-2\sqrt{\sqrt2-1}}$}
On pose $x+1=\sqrt2$, soit $x=\sqrt2-1\approx0{,}414\in[0;1]$. L'expression est alors $A$ pour cette valeur de $x$, et d'apr\`es la question d. :
\[
\sqrt{\sqrt2+2\sqrt{\sqrt2-1}}-\sqrt{\sqrt2-2\sqrt{\sqrt2-1}}=2\sqrt{\sqrt2-1}.
\]

\subsection*{f. Calcul de $A$ pour $x\ge1$, puis application}
Pour $x\ge1$, on a $\sqrt x\ge1$, donc $\lvert\sqrt x-1\rvert=\sqrt x-1$ et
\[
A=(\sqrt x+1)-(\sqrt x-1)=2.
\]
Pour $x=2024$ : $x+1=2025$ et $2\sqrt x=2\sqrt{2024}$, donc
\[
\sqrt{2025+2\sqrt{2024}}-\sqrt{2025-2\sqrt{2024}}=A(2024)=2.
\]

%% =====================================================================
\section*{Exercice 3}

\subsection*{1. Calculs}

\paragraph{a. $A$ sous forme de fraction irr\'eductible.}
\[
8+\frac73-\frac32=\frac{48+14-9}{6}=\frac{53}{6},\qquad
\frac15-3\cdot\frac92=\frac15-\frac{27}{2}=-\frac{133}{10}.
\]
\[
A=\frac{53}{6}\times\left(-\frac{133}{10}\right)=-\frac{7049}{60}.
\]
Comme $7049=7\times19\times53$ et $60=2^2\times3\times5$, la fraction est irr\'eductible.

\paragraph{b. $B$ simplifi\'e ($a\neq0$, $b\neq0$).}
\[
\frac{a^3b^{-2}}{a^4b^{-3}}=a^{-1}b,\qquad \frac{a^2b^3}{a^4b^{-3}}=a^{-2}b^6,
\]
donc
\[
B=a^{-3}b^{7}=\frac{b^7}{a^3}.
\]

\paragraph{c. $C$ en puissances de nombres premiers.}
On \'ecrit $0{,}6=\dfrac35$, $0{,}8=\dfrac{2^2}{5}$ et $0{,}4=\dfrac25$.
\[
\text{Num\'erateur : }(0{,}6)^2\times12^5\times54^3=\frac{3^2}{5^2}\times(2^2\cdot3)^5\times(2\cdot3^3)^3=\frac{2^{13}\cdot3^{16}}{5^2},
\]
\[
\text{D\'enominateur : }9^2\times5^3\times(0{,}8)^3\times(0{,}4)^4=3^4\cdot5^3\cdot\frac{2^6}{5^3}\cdot\frac{2^4}{5^4}=\frac{2^{10}\cdot3^4}{5^4}.
\]
\[
C=\frac{2^{13}\cdot3^{16}}{5^2}\times\frac{5^4}{2^{10}\cdot3^4}=2^3\times3^{12}\times5^2.
\]

\subsection*{2. \'Equations et in\'equations}

\paragraph{$(I_1)$ : $\lvert x+4\rvert\le2$.}
\[
-2\le x+4\le2\iff-6\le x\le-2.\qquad S_1=[-6\,;\,-2].
\]

\paragraph{$(I_2)$ : $\lvert x+6\rvert>2$.}
\[
x+6>2\ \text{ou}\ x+6<-2\iff x>-4\ \text{ou}\ x<-8.
\]
\[
S_2=\left]-\infty\,;\,-8\right[\ \cup\ \left]-4\,;\,+\infty\right[.
\]

\paragraph{(E) : $\lvert8-x\rvert=3$.}
\[
8-x=3\ \text{ou}\ 8-x=-3\iff x=5\ \text{ou}\ x=11.\qquad S_E=\{5\,;\,11\}.
\]

\subsection*{3. In\'egalit\'es avec $x,y,z\ge0$}

\paragraph{a.} $x+y-2\sqrt{xy}=(\sqrt x-\sqrt y)^2\ge0$, donc $x+y\ge2\sqrt{xy}$.

\paragraph{b.} On applique a) aux trois couples :
$x+y\ge2\sqrt{xy}$, $x+z\ge2\sqrt{xz}$, $y+z\ge2\sqrt{yz}$.
Ces quantit\'es sont positives, on peut donc les multiplier :
\[
(x+y)(x+z)(y+z)\ge8\sqrt{x^2y^2z^2}=8xyz.
\]

\paragraph{c.} Pour $x,y>0$ : $x^2+y^2\ge2xy>0$, donc $\dfrac{1}{x^2+y^2}\le\dfrac{1}{2xy}$.

\paragraph{d.} Comme $x+y\ge0$, on multiplie c) par $x+y$ :
\[
\frac{x+y}{x^2+y^2}\le\frac{x+y}{2xy}=\frac12\left(\frac1x+\frac1y\right).
\]

\subsection*{4. Encadrements et simplification}

On a $1\le a\le3$ et $-7\le b\le-1$.

\paragraph{a. Encadrements.}
\begin{itemize}
  \item $4-a$ : $1\le4-a\le3$.
  \item $2a-b$ : $2\le2a\le6$ et $1\le-b\le7$, donc $3\le2a-b\le13$.
  \item $ab-1$ : comme $b<0$, $ab\in[-21\,;\,-1]$, donc $-22\le ab-1\le-2$.
\end{itemize}

\paragraph{b. Simplification de $D$.}
$\sqrt{(4-a)^2}=4-a$ (car $4-a>0$), $\sqrt{(2a-b)^2}=2a-b$ (car $2a-b>0$), et $\lvert ab-1\rvert=1-ab$ (car $ab-1<0$). Donc
\[
D=(4-a)+(2a-b)-(1-ab)=3+a-b+ab.
\]

\subsection*{5. Racines et identit\'e}

On utilise $A=\sqrt{x+1+2\sqrt x}-\sqrt{x+1-2\sqrt x}$ (\'enonc\'e peu lisible sur la photo, d\'efinition reconstitu\'ee).

\paragraph{a. D\'eveloppement.} $(\sqrt x-1)^2=x-2\sqrt x+1$.

\paragraph{b. D\'eduction.} Comme $(\sqrt x-1)^2\ge0$, on a $x-2\sqrt x+1\ge0$, soit
\[
x+1\ge2\sqrt x.
\]

\end{document}
